% ---------------------------------------------------------------------------
% The transport assembly: the W-form, the ell=0 row, and the kappa' ledger.
%
% SELF-CONTAINED BY DESIGN. Inputs one file, the history figure, and defines
% nothing. Written 2026-10-21 by the Paper 1 stream at the PI's direction.
%
% REVISED 2026-10-22, same hand. Three statements in the 10-21 text were
% carried past their own refutations:
%   - T-105, the restored a.B_0 row -> WITHDRAWN. Refuted by T-107, by the
%     tau_0 = 0 constraint that Appendix (frame-division) already carried,
%     and by criterion 5b doubling at 98.5%. The ell=0 result is now derived
%     WITHOUT it and is cleaner: C_0 = -(1/3) a.B_2, one correction, T-104.
%   - the kappa' ledger as a PROHIBITION -> it is a REVIEW FLAG. T-113, S48.
%     W(kappa' S) = V S is exact, so a coefficient MAY hold kappa'.
%   - the by-parts slots as ARTEFACTS -> an exact by-parts preserves the
%     transfer. The slots are a legitimate step performed inexactly; c3
%     measures that at 10-250%.
% Plus two missing qualifiers: GDE99 (37) is first order, and kappa'' =
% -kappa'^2 holds at the visibility maximum, not generally.
%
% SUPERSESSIONS ARE NOW STATED AT THE POINT OF USE, not collected in this
% comment where no reader meets them. That is the master file's own pattern
% (its lines 347-354) and it is better than what this section did.
%
% FINDING NUMBERS USED, ALL EXISTING, NONE MINTED HERE: T-103, T-104, T-107,
% T-113, T-117, S48, B-11, B-16, T-101/D18, T-88, T-96, T-85, T-66, A-5.
% The bundle stream owns provenance/BUNDLE-DEFECTS-v1.md.
% ---------------------------------------------------------------------------
\section{The transport assembly}
\label{sec:transport-assembly}

\emph{The reader who wants only the result: \S\ref{sec:w-form} is the
assembly in which the by-parts does not have to be taken,
\S\ref{sec:ell-zero-row} corrects (90)'s $C_0$ by a factor $\tfrac{15}2$, and
\S\ref{sec:kappa-ledger} is the structural check that flags both before they
are implemented.}

Three questions about the line-of-sight solution are one question, and this
section answers it once. They are: which weight carries which source; what the
$\ell=0$ matching row of (89) should be; and why every repair attempted on
(102) so far has moved a weight without moving the term that dominates the
integrand. The common answer is that \textbf{(102)'s integration by parts was
taken onto the coefficients}, and that the assembly in which that step never
arises is the one obtained by multiplying the transport equation by its own
integrating factor.

% ---------------------------------------------------------------------------
% The photon's history, radiation era through decoupling to free streaming.
%
% SELF-CONTAINED BY DESIGN. Inputs nothing. Requires tikz with the libraries
% positioning, arrows.meta, fit, backgrounds, calc.  Written 2026-10-21 by the
% Paper 1 stream at the PI's direction; it is the spine the theory sections
% are ordered on, and every box names the section that treats it.
%
% GRAYSCALE-SAFE. Nothing is distinguished by colour alone: the three source
% weights are distinguished by line style (solid / dashed / dotted) and by
% their printed label, and every defect box carries its letter.
% ---------------------------------------------------------------------------
\begin{figure*}[tb]
\centering
\resizebox{\textwidth}{!}{%
\begin{tikzpicture}[
  font=\footnotesize,
  >={Stealth[length=5pt]},
  epoch/.style={draw, rounded corners=1pt, align=center, minimum height=7mm,
                inner sep=3pt, fill=black!4},
  phys/.style={align=center, inner sep=2pt},
  slot/.style={draw, align=center, inner sep=3pt, minimum height=6mm,
               font=\scriptsize},
  defect/.style={draw, dashed, align=left, inner sep=3pt, font=\scriptsize,
                 rounded corners=1pt},
  sec/.style={font=\scriptsize\itshape, text=black!55},
  tick/.style={draw, black!55},
]

% ---- the conformal-time spine --------------------------------------------
\def\XA{0}  \def\XB{3.0} \def\XC{6.2} \def\XD{8.8} \def\XE{12.0} \def\XF{15.4}
\draw[->, thick] (\XA-0.5,0) -- (\XF+0.9,0);
\node[anchor=west] at (\XF+1.0,0) {$\eta$};

\foreach \x/\lab in {\XB/{$\eta_{\rm eq}$}, \XD/{$\eta_*$},
                     \XE/{$6\eta_*$}, \XF/{$\eta_0$}} {
  \draw[tick] (\x,-0.13) -- (\x,0.13);
  \node[below=1pt, font=\scriptsize] at (\x,-0.13) {\lab};
}

% ---- epoch band -----------------------------------------------------------
\node[epoch, anchor=south west] at (\XA-0.5,0.30)
  {radiation era};
\node[epoch, anchor=south west] at (\XB+0.08,0.30)
  {tight coupling};
\node[epoch, anchor=south west] at (\XD+0.08,0.30)
  {slow decoupling};
\node[epoch, anchor=south west] at (\XE+0.08,0.30)
  {free streaming};

% ---- what the hierarchy is, stage by stage --------------------------------
\node[phys, anchor=south west] at (\XA-0.5,1.15)
  {\begin{minipage}{3.2cm}\centering
   $\Phi$ decays through\\horizon entry;\\the oscillator is \emph{driven}
   \end{minipage}};
\node[phys, anchor=south west] at (\XB+0.05,1.15)
  {\begin{minipage}{2.8cm}\centering
   $\tau_{A_\ell}\!\approx\!0$, $\ell>2$;\\
   $\tau_{ab}\simeq-\kappa'^{-1}\!\big[\sigma_{ab}\!+\!\mathrm D_{\langle a}v^B_{b\rangle}\big]$
   \end{minipage}};
\node[phys, anchor=south west] at (\XD+0.05,1.15)
  {\begin{minipage}{2.9cm}\centering
   the visibility $\mathcal V$\\has finite width\\$\Delta\eta=0.00972$
   \end{minipage}};
\node[phys, anchor=south west] at (\XE+0.05,1.15)
  {\begin{minipage}{3.0cm}\centering
   $\kappa'\!\to\!0$, $W\!\to\!1$;\\GDE99~(37) holds\\
   \textsc{theorem}
   \end{minipage}};

% ---- the two weights, drawn as what they are ------------------------------
\node[anchor=east, font=\scriptsize] at (\XA-0.6,-1.15) {$W=e^{-\kappa}$};
\draw[thick, domain=\XA-0.5:\XF, samples=160, smooth]
  plot (\x, {-1.15 + 0.55/(1+exp(max(-8,min(8,-(\x-\XD)*2.2))))});
\node[anchor=east, font=\scriptsize] at (\XA-0.6,-2.35) {$\mathcal V=W\kappa'$};
\draw[thick, dashed, domain=\XA-0.5:\XF, samples=200, smooth]
  plot (\x, {-2.35 + 0.95*exp(max(-8,-((\x-\XD)*1.55)^2))});
\draw[tick, dotted] (\XD,-2.45) -- (\XD,0);

% ---- the W-form: which source travels under which weight -------------------
\node[slot, anchor=north west, text width=5.9cm, align=left] at (\XA-0.5,-3.15)
  {\textbf{under $W$ --- the gravitational source $a\mathcal B_\ell$}\\[1pt]
   $a\mathcal B_0=-\tfrac13\mathrm D^a\tau_a$,\quad
   $a\mathcal B_1=\mathrm D_a\ln T+A_a$,\\
   $a\mathcal B_2=\sigma_{ab}$\quad(all three vanish in the background)};
\node[slot, anchor=north east, text width=5.9cm, dashed, align=left]
  at (\XF+0.9,-3.15)
  {\textbf{under $\mathcal V$ --- the in-scattered source $\mathcal S_\ell$}\\[1pt]
   $\mathcal S_1=\bar f_1v^B_a$,\quad
   $\mathcal S_2=\bar f_2\tau_{ab}$\\
   \emph{Annals~II} sets $\bar f_2=0$: that is A-5, and it is \textsc{triage c}
   --- $\Delta\mathcal S_2=0$ by GDE99~(37)};
\node[anchor=north, align=center, font=\footnotesize]
  at ({(\XA+\XF)/2}, -4.62)
  {$\dfrac{\mathrm d}{\mathrm d\eta}\big[W\tau_\ell\big]
    + W\,(\text{streaming})
    \;=\; W\,a\mathcal B_\ell \;+\; \mathcal V\,\mathcal S_\ell$};
\node[anchor=north, align=center, font=\scriptsize]
  at ({(\XA+\XF)/2}, -5.25)
  {\textbf{One $\kappa'$, consumed by the integrating factor. No coefficient
   holds one, so none can be differentiated.}};

% ---- the defects, each at the conformal time it acts at --------------------
\node[defect, anchor=north west, text width=2.9cm] at (\XA-0.5,-6.05)
  {\textbf{C3} radiation driving\\
   prescribed $\Phi$, static $T(k)$\\
   \emph{amplitude and phase}\\\textsc{open} (A-3)};
\node[defect, anchor=north west, text width=2.9cm] at (\XB+0.45,-6.05)
  {\textbf{C5} $k_D$ taken at $R=0$\\
   while $R(\eta_*)\simeq0.61$\\
   \emph{the damping scale}\\\textsc{measured}, T-81};
\node[defect, anchor=north west, text width=3.0cm] at (\XD-0.95,-6.05)
  {\textbf{A4\,/\,A5} the $\ell=0$ row\\
   $\alpha_2$ where $\alpha_0$ belongs;\\
   $a\mathcal B_0\delta_{\ell0}$ dropped\\\textsc{theorem}, \texttt{c5}};
\node[defect, anchor=north west, text width=2.9cm] at (\XE-0.55,-6.05)
  {\textbf{C4} frozen at $6\eta_*$\\
   no slip, no quadrupole,\\no polarisation\\\textsc{open} (A-5)};
\node[defect, anchor=north west, text width=2.9cm] at (\XF-0.1,-6.05)
  {\textbf{C6} $\mathcal D(k)$ applied\\
   where $e^{-\kappa}=1$\\
   \emph{the late ISW}\\\textsc{open}, T-81b};

% the one defect that is not local in eta
\node[defect, anchor=north, text width=13.4cm, thick]
  at ({(\XA+\XF)/2+1.0}, -7.95)
  {\textbf{A1 --- not local in $\eta$, and the one that is upstream of the rest.}
   (102) carries $\kappa'$ inside $C_1$ and then differentiates it, so the
   released integrand holds $\mathcal V\kappa''\hat v_B$ at $93$--$97\%$ of its
   own total: \textbf{three powers of $\kappa'$ where the source carries one}.
   The $W$-form above is the assembly in which that term cannot arise, because
   no coefficient holds $\kappa'$ to be differentiated. Instrument:
   \texttt{tests/frame/c7\_kappa\_ledger.py}.};

% ---- section index --------------------------------------------------------
\node[sec, anchor=north west] at (\XA-0.5,-9.55)
  {§\,\ref{sec:w-form} the $W$-form\ \ \textbar\ \
   §\,\ref{sec:ell-zero-row} the $\ell=0$ row\ \ \textbar\ \
   §\,\ref{sec:kappa-ledger} the $\kappa'$ ledger\ \ \textbar\ \
   §\,\ref{sec:frame-discrepancy} triage\ \ \textbar\ \
   §\,\ref{sec:scattering-expansion} the scattering expansion};

\end{tikzpicture}%
}
\caption{The photon's history and where each open question acts.$^{\dagger}$}
\label{fig:history-flow}
\end{figure*}


Figure~\ref{fig:history-flow} is the order this section and
\S\ref{sec:frame-discrepancy} are written in. It is the photon's history and
not a table of contents: each open question is drawn at the conformal time it
acts at, which is what separates the questions that a frame instrument can see
from the questions it cannot.

% ===========================================================================
\subsection{The $W$-form}
\label{sec:w-form}

\paragraph{The derivation, in three lines.} The covariant transport equation
at rank $\ell$ is
\begin{equation}
  \tau_\ell' \;+\; (\text{streaming couplings}) \;+\; \kappa'\tau_\ell
  \;=\; a\mathcal B_\ell \;+\; \kappa'\mathcal S_\ell ,
  \label{eq:transport}
\end{equation}
with $a\mathcal B_\ell$ the gravitational source of
\S\emph{The generic source derivatives}' rank decomposition and
$\mathcal S_\ell$ the in-scattered Thomson source. Both
$\kappa'$ on display are the same $\kappa'$: the out-scattering on the left and
the in-scattering on the right are two halves of one collision term.

Multiply through by $W=e^{-\kappa}$. With $\kappa$ accumulated backwards from
$\eta_0$, so that $\mathrm d\kappa/\mathrm d\eta=-\kappa'$ and
$W'=+\kappa'W=\mathcal V$ --- the sign this bundle's convention forces, stated
at \S\ref{sec:frame-discrepancy} --- the first and third terms of
\eqref{eq:transport} combine exactly:
\begin{equation}
  W\tau_\ell' + W\kappa'\tau_\ell \;=\; W\tau_\ell' + W'\tau_\ell
  \;=\; \frac{\mathrm d}{\mathrm d\eta}\big[W\tau_\ell\big] .
\end{equation}
Hence
\begin{equation}
  \boxed{\;
  \frac{\mathrm d}{\mathrm d\eta}\big[W\tau_\ell\big]
  \;+\; W\,(\text{streaming couplings})
  \;=\; W\,a\mathcal B_\ell \;+\; \mathcal V\,\mathcal S_\ell \;}
  \label{eq:w-form}
\end{equation}
and integrating from $\eta_*$ to $\eta_0$ gives the solution with
$\big[W\tau_\ell\big]_{\eta_*}^{\eta_0}
=\tau_\ell(\eta_0)-W(\eta_*)\tau_\ell(\eta_*)$ as its only boundary term,
since $W(\eta_0)=1$.

\paragraph{What \eqref{eq:w-form} says, and it is three things.}

\begin{enumerate}\itemsep3pt
\item \textbf{The gravitational source travels under $W$ and the in-scattered
      source under $\mathcal V$.} Neither weight is assigned; both fall out of
      one multiplication. $\mathcal V$ appears because $\kappa'$ was already
      beside $\mathcal S_\ell$ in the collision term, and $W\kappa'=\mathcal V$
      is an identity.
\item \textbf{$\kappa'$ is consumed once}, by the integrating factor, and
      appears in no coefficient. So no differentiation of a coefficient can
      produce a $\kappa'^2$ --- which is the error class
      \S\ref{sec:kappa-ledger} flags, not a theorem about where $\kappa'$ may
      live.
\item \textbf{The boundary term is one object, and it is exponentially
      suppressed at $\eta_*$.} $\big[W\tau_\ell\big]$, not
      $\big[C_2'-C_1\big]$ plus $kC_2$, and it carries $W$ rather than a
      diffusion factor. \textbf{D18} measures what the released boundary
      carries instead: with \texttt{boundary\_weight\_w} \texttt{False},
      $C_1(\eta_*)$ enters under the diffusion factor alone --- no $W$, no
      $\mathcal V$ --- so its bare $\kappa'(\eta_*)$, $3.82\times10^2$ on
      Einstein--de~Sitter and $2.64\times10^2$ on \texttt{PLANCK\_MODEL},
      stands under no exponential suppression at all. \eqref{eq:w-form} has no
      counterpart: $\mathcal V\mathcal S_1$ is $\kappa'$-free.
\end{enumerate}

\paragraph{What \eqref{eq:w-form} does \emph{not} say, corrected here.}
An earlier draft of this section called (102)'s three bulk slots and boundary
pair \emph{artefacts of a step \eqref{eq:w-form} does not take}. \textbf{That
overstates and is withdrawn.} An \emph{exact} integration by parts preserves
the transfer function, so the slots are not artefacts: they are the correct
consequence of a legitimate step. The defensible claim is narrower and is the
one that matters --- \textbf{(102) as printed is not the exact by-parts of its
own integral}, which \texttt{c3\_byparts.py} measures at $10$ to $250\%$ in
the printed bulk. \eqref{eq:w-form} is therefore not a different physics; it is
a construction in which that arithmetic cannot go wrong, because there are no
coefficients to take the by-parts onto.

\paragraph{This supersedes \S\ref{sec:b-slot}'s item~2, and the wording that
item carries is what \textbf{T-88} implemented.} That item reads \emph{``they
take different weights: $W$ for the gravitational half, $g=\mathcal V$ for the
scattering half''}. As a description of the outcome of \eqref{eq:w-form} it is
correct. As an instruction to a solver it is not, and it was read as one:
multiplying $C_1^{\rm thom}=-\kappa'\hat v_B$ by $\mathcal V$ applies a second
$\kappa'$, which \S\ref{sec:repair-attempts} measures at $320$ to $682$.
\textbf{The scattering half arrives at $\mathcal V$ weight carried on the
undifferentiated $\hat v_B$; it is not multiplied by $\mathcal V$.} Finding
\textbf{T-103}; the phrasing is the Paper~1 stream's and the correction is
recorded against it rather than against the implementation that followed it.

\paragraph{What $\mathcal S_2$ is, and why it cannot be a frame defect.}
\emph{Annals~II}'s collision term is
$C[\tau]\simeq\sigma_{\rm T}n_E\big(\bar f_1e^av^B_a
+\bar f_2e^ae^b\tau_{ab}-\tau\big)$ with $\bar f_1=1$ and
$\bar f_2=0$, so $\mathcal S_2=0$ in this reconstruction. That is
approximation \textbf{A-5} stated at the collision term rather than at the
oscillator. \textbf{It cannot produce a frame disagreement at linear order}:
$\mathcal S_2\propto\tau_{ab}$, and GDE99~(37) leaves $\tau_{A_\ell}$ invariant
for $\ell\ge2$ \emph{at first order in $v$} --- those multipoles are already
first-order small, so what a boost would generate from them is second order ---
whence $\Delta\mathcal S_2=0$ and dropping the term changes both columns by the
same amount. It is a \textbf{\S\ref{sec:frame-discrepancy} Triage~C} site ---
it can ruin the spectrum and it cannot be seen by a frame instrument --- and
listing it as a candidate for the frame gap confuses the two questions the
triage tables exist to separate.

% ===========================================================================
\subsection{The $\ell=0$ matching row}
\label{sec:ell-zero-row}

\paragraph{The rows of (89) are Taylor coefficients of the kernel at the
origin, and one of the three uses the wrong rank's normalisation.}
With $h_\ell=\alpha_\ell j_\ell$ and
$\alpha_\ell=(2\ell+1)/\beta_\ell$ --- this bundle's own
\texttt{functions/harmonics/weights.py}, where
$\beta_\ell=\ell!/(2\ell-1)!!$ gives $\beta_0=\beta_1=1$ and
$\beta_2=\tfrac23$ --- each row of (89) pairs $C_n$ with
$h_\ell^{(n)}(0)$:

\begin{center}\small
\begin{tabular}{@{}cccl@{}}
\toprule
row & used & its own & the entry \\
\midrule
$\ell=2$ & $\alpha_2=\tfrac{15}2$ & $\alpha_2$ &
  $h_2''(0)=\tfrac{15}2\cdot\tfrac2{15}=1$ \\
$\ell=1$ & $\alpha_1=3$ & $\alpha_1$ &
  $h_1'(0)=3\cdot\tfrac13=1$ \\
$\ell=0$ & $\boldsymbol{\alpha_2=\tfrac{15}2}$ & $\boldsymbol{\alpha_0=1}$ &
  $-\tfrac52$ printed; $h_0''(0)=-\tfrac13$ \\
\bottomrule
\end{tabular}
\end{center}

\noindent
(89) writes its $\delta_{\ell0}$ entry as $5\beta_2^{-1}\cdot\tfrac13$, which
is $\alpha_2$ times $j_0''(0)$ up to the sign \textbf{B-11} corrects. The
rank-two normalisation is right where it belongs and is carried into a row
that is not rank two. \textbf{The two rows that are exactly unity are the
control}, and they are a control precisely because they come out at unity
\emph{without adjustment}: a normalisation error cannot produce $1$ twice and
then $-\tfrac52$ once. Findings \textbf{B-16} and \textbf{T-104}; the
discrepancy is $\alpha_2/\alpha_0=\tfrac{15}2$.

\textbf{This route is independent of (89), which is what makes it a check
rather than a restatement.} $\beta_\ell$ is fixed by the PSTF normalisation
$O^{A_\ell}O_{A_\ell}$ and $j_0''(0)=-\tfrac13$ by the ascending series;
neither comes from (89). (89) is the \emph{claim} about what the rows hold and
the kernel is the \emph{object}. Two of the three rows reproduce (89)'s entry
exactly from the object, which validates the route; the third does not.

\paragraph{The right-hand side of the $\ell=0$ row is zero, and that is forced
by a definition.} $\tau_0\equiv0$ is a \textbf{constraint, not initial data}:
$T(x)$ is defined as the monopole of $T(x,e)$, so differentiating
$\int A_0\,\mathrm d\eta'=0$ puts zero on the right of the row.
$\mathcal B_0$ is a functional of the solution, not an external forcing.
Appendix~\ref{sec:frame-division} states this and states it first ---
\emph{``(89)'s $\ell=0$ line is not mis-sourced \ldots $C_0$ is correctly
derived''}. With $\alpha_0$ on the left and zero on the right,
\begin{equation}
  C_0 + k^2C_2\,h_0''(0) = 0 ,
  \qquad h_0''(0)=-\tfrac13 ,
\end{equation}
and with $C_2=-k^{-2}a\mathcal B_2$ this solves to
\begin{equation}
  \boxed{\;C_0 \;=\; -\tfrac13\,a\mathcal B_2\;}
  \label{eq:c0-solved}
\end{equation}
in place of (90)'s $C_0=-\tfrac52a\mathcal B_2$. \textbf{One correction, and
the ratio against the printed row is $\tfrac52\big/\tfrac13=\tfrac{15}2$
exactly} --- $\alpha_2/\alpha_0$, with no $\mathcal B_0$ anywhere.
Instrument \texttt{c5\_monopole.py}, under \texttt{tests/frame/}, which carries
a negative control on each alternative.

\paragraph{Where $\mathcal B_0$ does enter, and why restoring it to $C_0$ is
double counting. \textnormal{\textbf{Correction, 2026-10-22.}}} An earlier
draft of this subsection restored an $a\mathcal B_0\delta_{\ell0}$ row,
printed $C_0=-\tfrac13a\mathcal B_2-a\mathcal B_0$, and recorded it as
\textbf{T-105}. \textbf{That is withdrawn.} By identity
\begin{equation}
  a\mathcal B_0=-\tfrac{k}{3}\,\tau_1 ,
\end{equation}
derived in index algebra and controlled on rank~1, where the identical chain
returns $a\mathcal B_1=k(\delta T+\Phi_A)$ --- the released
\texttt{coefficients.py} line. And (107) carries
$-C_1'=(\kappa'\hat v_B)'+(a\mathcal B_1)'/k$ with
$(a\mathcal B_1)'/k=(\Phi_A'-\Phi_H')+\tfrac k3\tau_1$, so
\textbf{the monopole source is already inside the bracket, in $-C_1'$.}
Adding it to $C_0$ counts it twice, and criterion 5b doubles at $98.5\%$.
Finding \textbf{T-107} replaces T-105.

\paragraph{So $\Delta C_0\neq0$, and that is the correct outcome.} Under
\eqref{eq:c0-solved}, with Appendix~\ref{sec:frame-division}'s own
$\Delta(a\mathcal B_2)=-kv$ and $\Delta C_2=v/k$,
\begin{equation}
  \Delta C_0=-\tfrac13\Delta(a\mathcal B_2)=+\tfrac{kv}{3},
  \qquad
  \gamma\equiv\frac{\Delta C_0}{k^2\Delta C_2}=\tfrac13 .
\end{equation}
\textbf{A non-zero $\gamma$ is the signature that $\Delta C_0$'s partner is
cross-coefficient}, sitting in $-C_1'$ rather than in $C_2$. A $\gamma$ of
zero would say the partner was in $C_2$, which it is not. The withdrawn draft
reached $\gamma=0$ by \emph{adding a term}, and a gate that can be driven to
zero by adding a term is a fit rather than a gate --- which is \textbf{M15},
\emph{a gate built from one coefficient cannot see a cancellation distributed
across the bracket}. \texttt{c8\_rank0\_partner.py} scores the whole rank-0
content of the bracket and is the gate for this row;
Appendix~\ref{sec:frame-division}'s questions~3 and~4 close with \textbf{c8} or
not at all. Its questions~1 and~2 are settled by \textbf{T-66} and by
\S\ref{sec:w-form} respectively.

\paragraph{One thing \eqref{eq:c0-solved} does not disturb.} Appendix
\ref{sec:frame-division} excludes a comoving/physical or amplitude/potential
mismatch by measuring $C_0/(k^2C_2)$ as a background-independent and
$k$-independent rational. That argument holds under
\eqref{eq:c0-solved} unchanged --- the ratio is still a pure rational, now
$\tfrac13$ rather than $\tfrac52$ --- because what it rests on is that $C_0$
and $C_2$ divide one $a\mathcal B_2$. \textbf{The measurement was right and
what it established was internal consistency, which is not correctness.} It is
the same reading the bundle gave T-82 and T-85, applied to itself.

% ===========================================================================
\subsection{The $\kappa'$ ledger}
\label{sec:kappa-ledger}

\paragraph{The flag, and it is a flag and not a theorem.
\textnormal{\textbf{Revised 2026-10-21, T-113 and S48.}}} \emph{A coefficient
carrying a bare $\kappa'$ is a site where a large residual is manufactured the
moment it is differentiated, and is worth flagging before implementation.} An
earlier draft of this subsection printed that as a prohibition --- \emph{``$\kappa'$
may live in $W$, in $\mathcal V$, or in a derivative of either, never in a
coefficient''} --- and it was read as normative. \textbf{It is not a theorem},
because
\begin{equation}
  W\big(\kappa'\mathcal S\big)=\mathcal V\,\mathcal S
  \label{eq:ledger-identity}
\end{equation}
is exact: a coefficient \emph{may} hold $\kappa'$, and
$(\kappa'\mathcal S)'=\kappa''\mathcal S+\kappa'\mathcal S'$ is not
intrinsically spurious. The defect in (102) is wrong weights, signs,
boundaries or missing partners, not the presence of $\kappa'$.

\paragraph{Why the flag is kept.} It separates two $\kappa'^2$ that are
indistinguishable in the output. The one inside
$\dot{\mathcal V}=(\kappa''+\kappa'^2)W$ comes from differentiating the
\textbf{weight}; it is required, and \eqref{eq:w-form} produces it whenever a
streaming term is integrated. The one from differentiating a
\textbf{coefficient} is not required by anything. Only the provenance
distinguishes them, so the flag is raised on the provenance. At $\eta_*$ in
this bundle's units \textbf{D18} measures $\kappa'=3.82\times10^2$ on
Einstein--de~Sitter and $2.64\times10^2$ on \texttt{PLANCK\_MODEL}, so a
spurious power of it is two to three orders. \textbf{What the flag may no
longer do is settle a question on its own.}

\paragraph{(102) as printed trips it, and \textbf{T-117} is the review.}
\S\emph{The generic source derivatives} gives
$-C_1\simeq k^{-1}(a\mathcal B_1+\kappa'v_B)$ --- a coefficient holding a bare
$\kappa'$ --- and (102) then differentiates it.
\S\ref{sec:frame-discrepancy} measures the consequence:
$\mathcal V\kappa''\hat v_B$ is $93$--$97\%$ of the released integrand. The
flag asks whether that is a defect; the super-horizon limit answers it where
the size is decidable rather than arguable. In the Newtonian column as
$k\to0$, $\kappa'v_B/k$ \textbf{does not die} --- it tends to $-8.88$, which is
$-26.6\times$ the Sachs--Wolfe term --- and the bulk is then entirely
$\int\mathcal V(\kappa'\hat v_B)'j_\ell$ at $9.46\times10^3$.
\textbf{One object, both terms.} The flag was raised on the right site; what it
could not supply, and \textbf{T-117} does, is the size.

\paragraph{A local statement, not a general one.} $\kappa''=-\kappa'^2$ holds
\textbf{at the visibility maximum}, where $\dot{\mathcal V}=0$. So the reading
of $\mathcal V\kappa''\hat v_B$ as \emph{three powers of $\kappa'$} is local to
that point, and the measured $93$--$97\%$ share follows because $\mathcal V$
concentrates the integrand there. It is not a statement about the integrand at
arbitrary $\eta$.

\paragraph{It is applied before implementation, not after.}
\texttt{tests/frame/c7\_kappa\_ledger.py} grades an assembly structurally ---
for each term, whether its coefficient holds $\kappa'$ and at what power ---
rather than measuring an output. It reports \eqref{eq:w-form} clean, and (102)
as printed and T-88's variant as carrying $\kappa'$ in a coefficient at power
two. \textbf{Both of the defects of 19--20~October would have been flagged by
that one check before either was implemented}, which is the argument for a
structural review alongside a numerical one: a numerical check on an assembly
that is wrong in this way measures the defect faithfully and reports a number.

% ===========================================================================
\subsection{What this section establishes, and what it does not}

\textbf{Established.} That \eqref{eq:w-form} follows from
\eqref{eq:transport} by one multiplication; that it assigns both weights
without a choice and leaves no $\kappa'$ in any coefficient; that (89)'s
$\ell=0$ row carries $\alpha_2$ where $\alpha_0$ belongs, with the $\ell=1$ and
$\ell=2$ rows as unadjusted controls; that the row's right-hand side is zero by
the $\tau_0\equiv0$ constraint, giving \eqref{eq:c0-solved}; that
$a\mathcal B_0=-\tfrac k3\tau_1$ is already carried in $-C_1'$, so restoring it
to $C_0$ double counts; and that $\mathcal S_2$ is a Triage~C site and cannot
carry a frame gap at linear order.

\textbf{Not established.} That \eqref{eq:w-form} closes the frame gap. It
satisfies the ledger by construction and no setting of
\texttt{integral\_solution}'s seven controls reaches it --- each redistributes
weights while leaving $\kappa'$ inside $C_1$ --- but that is an argument from
exhaustion of what has been tried, not a measurement. \textbf{And it is now
known not to be sufficient}: \textbf{T-117} finds the two frames failing by
different mechanisms in the super-horizon limit, the Newtonian column through
the $\kappa'$ coefficient and the CDM column through a sign in the boundary at
$-0.9999$ against a target of $+1$. One repair does not obviously reach both.

\textbf{Also not established, and it is a dependency of
\eqref{eq:c0-solved}.} \eqref{eq:c0-solved} \emph{uses}
$C_2=-k^{-2}a\mathcal B_2$, which comes from the same solve in which
\textbf{B-11} found a sign error and \textbf{T-104} a normalisation error. The
$\ell=2$ row returning $h_2''(0)=1$ controls the normalisation and not the
sign or the source pairing. \textbf{No instrument yet audits the $\ell=2$ row
the way \texttt{c5} audits $\ell=0$}, and until one does, a defect of that
class in $C_2$ would be inherited silently here.

\textbf{The acceptance condition, stated in advance so the repair can fail
visibly.} (i)~$A_N=A_C$ to the quadrature floor \textbf{T-85} measures, neither
better nor worse; (ii)~the Newtonian column moves by no more than that floor;
(iii)~\texttt{c7\_kappa\_ledger.py} reports no coefficient holding $\kappa'$
--- \emph{a review signal that the error class is absent by construction, not a
proof of correctness}, since \eqref{eq:ledger-identity} makes the ledger a flag;
(iv)~\texttt{c8\_rank0\_partner.py}, not $\gamma$, scores the rank-0 content of
the bracket on the live coefficients; and (v)~\textbf{the spectrum adjudicates,
and it adjudicates last}, since two wrong columns can agree.

\textbf{The order is forced, not preferred --- in implementation, not in
exposition.} \S\ref{sec:ell-zero-row} is implemented before
\S\ref{sec:w-form}: the $\ell=0$ row sits inside the bracket the $W$-form
reweights, so a weight correction applied over a $C_0$ that is wrong by
$\tfrac{15}2$ measures the wrong thing. \textbf{T-96}'s second run is that
measurement --- a $W$-form instrument run against coefficients still carrying
the printed row, returning $O(1)$ --- and it is reported there as not
conclusive, correctly.
